\section{Which pattern emerges}
\label{sec:selection}

Every recipe is present in a multiplicative picture, and a bent front end
puts most of them into an additive one. Which of them do you see? Two
ingredients decide it. One is slowness: how small the recipe's heterodyne
ratio is. The other is strength: how much of the picture the recipe
carries, which for anything with edges falls off with the harmonics. The
theory of two sinusoids --- the textbook theory of beats --- has only the
first ingredient, because a sinusoid has one harmonic and nothing to
choose between; families with edges have both, and the visible pattern is
the one that is slow \emph{per unit of strength}.

\begin{definition}[The priced merit]
\label{def:merit}
For a recipe $(a,b)$ with $ab\ne0$ on two families with local rates
$\nabla\cnt_1$ and $\nabla\cnt_2$, the \emph{priced merit} is the heterodyne
ratio of Definition~\ref{def:eta} charged the product of the harmonic
orders that carry the recipe,
\[
  \mu_{a,b}\;=\;\het_{a,b}\cdot|ab|
  \;=\;\frac{\lvert a\,\nabla\cnt_1+b\,\nabla\cnt_2\rvert}
            {\tfrac12\big(\lvert a\,\nabla\cnt_1\rvert+\lvert b\,\nabla\cnt_2\rvert\big)}\cdot|ab| .
\]
The visible recipe is the one with the smallest priced merit, and it shows
as a beat when $\mu<\tfrac14$. For a first-order recipe $\mu=\het$, and the
rule is the four members per band of \S\ref{sec:slow}; a recipe of order
$ab$ is carried by harmonics $ab$ times weaker, and must be $ab$ times
slower to show as strongly.
\end{definition}

The price is Theorem~\ref{thm:whosees}'s second claim read backwards: the
recipe $(a,b)$ is carried by the $a$-th harmonic of one family and the
$b$-th of the other, each about $1/n$ strong, so the beat it makes has
about $1/|ab|$ of the contrast of the first-order beat. Without the price
the question has no answer. Every real ratio of pitches has rational
approximations as good as you like, so at any point some recipe of
enormous order is slower than every other, carried by harmonics no stroke
has --- a pair at the ratio $\sqrt2$ has a recipe $(41,-29)$ with more than
three thousand members per band and nothing to draw them with; with the
price, comparable slowness resolves toward the lower order, and the minimum
is a definite recipe at every point.

\begin{proposition}[Selection is best approximation]
\label{prop:selection}
For two parallel families of pitches $s_1$ and $s_2$, the recipes that beat
every lower-order recipe as the order is allowed to grow are exactly the
convergents of the continued fraction of $s_1/s_2$ (\selRecordsMatch\ of
\selRatios\ ratios measured, without exception). In the plane, where the
rates are vectors, the minimiser is found at every point by
Lagrange--Gauss reduction of the lattice $\{a\nabla\cnt_1+b\nabla\cnt_2\}$
--- the two-dimensional Euclidean algorithm, which finds a lattice's
shortest vectors in a handful of steps~\cite{NguyenStehle2009Gauss} ---
and the reduction names the brute-force winner in $99.8\%$ of
\selTrials\ random pairs (script \texttt{convergents.mjs}).
\end{proposition}

\begin{lookup}{continued fractions}
What to look up: \emph{continued fractions and convergents}, one page of
any number theory text. The facts used: a real $x$ has an expansion
$x=[a_0;a_1,a_2,\dots]$ into whole numbers, its \emph{partial quotients};
truncating gives the fractions $h_n/k_n$, the \emph{convergents} (here
$k_n$ is a denominator, not a recipe), which are the best rational
approximations of $x$ with denominators up to $k_n$; and the error of a
convergent is $|h_n-k_nx|=1/(x_{n+1}k_n+k_{n-1})$, where
$x_{n+1}=[a_{n+1};a_{n+2},\dots]$ is the \emph{complete quotient},
everything after the $n$-th term. For the golden ratio every partial
quotient is $1$; for $\sqrt2$, and for the silver ratio $1+\sqrt2$, every
one after the first is $2$; for $\pi$ they are $3;7,15,1,292,\dots$
\end{lookup}

The proposition is Lagrange's theorem on best approximation with a
different price: the classical statement prices a fraction $h/k$ by $k$,
the priced merit charges $hk$, and the records turn out to be the same
fractions. The reader who takes the one-dimensional case on trust loses
nothing that follows.

\subsection{Duty nulls}
\label{sec:dutynull}

Because a multiplicative picture's strength at a recipe is a product of
harmonics, a station is carried by exactly the harmonics its integers name,
and a family that lacks a harmonic cannot take part in a station that needs
it. That is a prediction the theory of sinusoids cannot make.

\begin{theorem}[Duty nulls]
\label{thm:dutynull}
Let the coarse family's pitch be $a/b$ times the fine family's, in lowest
terms. The visible station is the recipe $(a,-b)$: $a$ coarse counts minus
$b$ fine counts, carried by the coarse family's $a$-th harmonic and the fine
family's $b$-th. A family whose members cover a fraction $d$ of their pitch
has no $a$-th harmonic exactly when $ad$ is a whole number, so the station
is extinguished at those duties of the coarse family and at no others. A
$2{:}1$ pair (coarse pitch twice the fine) has $a=2$, $b=1$: its station
$(2,-1)$ dies at coarse duty $\tfrac12$; a $3{:}1$ pair's station $(3,-1)$
dies at $\tfrac13$ and $\tfrac23$.
\end{theorem}

\begin{proof}
With pitches $s_{\mathrm c}\approx(a/b)\,s_{\mathrm f}$ the combined count
$a\cnt_{\mathrm c}-b\cnt_{\mathrm f}$ has rate
$a/s_{\mathrm c}-b/s_{\mathrm f}\approx b/s_{\mathrm f}-b/s_{\mathrm f}=0$:
it is the slow one, and no recipe of lower order is. By
Theorem~\ref{thm:whosees} its strength is
$\hat g_{\mathrm c}(a)\,\hat g_{\mathrm f}(-b)$, and by the lookup of
\S\ref{sec:whosees}, $\hat g(n)=\sin(n\pi d)/n\pi$ vanishes exactly when
$nd$ is a whole number. A symmetric softening of the edges multiplies each
harmonic by a real factor and moves no zero.
\end{proof}

\begin{figure}[t]
  \centering
  \begin{tikzpicture}
      \begin{axis}[width=0.62\linewidth, height=4.9cm, xlabel={coarse duty $d$},
        ylabel={station strength}, xmin=0.18, xmax=0.82, ymin=0, ymax=0.045,
        scaled y ticks=false, ytick={0,0.01,0.02,0.03,0.04},
        tick align=outside, tick pos=left, axis line style={gray!60},
        grid=major, grid style={gray!18, very thin},
        label style={font=\footnotesize}, tick label style={font=\scriptsize},
        legend style={font=\scriptsize, draw=gray!40, at={(0.5,1.02)}, anchor=south, legend columns=2},
        xtick={0.25,0.333,0.5,0.667,0.75}, xticklabels={$\frac14$,$\frac13$,$\frac12$,$\frac23$,$\frac34$}]
        \addplot[accent, line width=0.7pt] table[x=duty, y=two, col sep=comma]{data/dutynull-law.csv};
        \addlegendentry{$2{:}1$, recipe $(2,-1)$}
        \addplot[cool, line width=0.7pt] table[x=duty, y=three, col sep=comma]{data/dutynull-law.csv};
        \addlegendentry{$3{:}1$, recipe $(3,-1)$}
        \addplot[accent, only marks, mark=*, mark size=1.3pt] table[x=duty, y=two, col sep=comma]{data/dutynull-curve.csv};
        \addplot[cool, only marks, mark=square*, mark size=1.1pt] table[x=duty, y=three, col sep=comma]{data/dutynull-curve.csv};
      \end{axis}
    \end{tikzpicture}\\[4pt]
  
\includegraphics[width=\linewidth]{dutynull-triple.png}
  \caption{\figlead{Duty nulls.} Above: the strength of the station against
  the coarse family's duty, measured on drawn strokes (marks) over the law
  $|\sin(a\pi d)/a\pi|\,|\hat g_{\text{fine}}(1)|$ (lines): the $2{:}1$
  station dies at $d=\tfrac12$, the $3{:}1$ station at $\tfrac13$ and
  $\tfrac23$, and the softened edge of a drawn stroke is the only departure.
  Below: what pooling sees of the $2{:}1$ pair at coarse duty $0.35$, $0.5$
  and $0.65$, expanded sixfold about its mean: present, extinguished,
  present. The nulls are $\dutyNullDepthTwo\times$,
  $\dutyNullDepthThreeLow\times$ and $\dutyNullDepthThreeHigh\times$ deep.
  Scripts \texttt{dutynull.mjs}, \texttt{dutynull-figs.mjs}.}
  \label{fig:dutynull}
\end{figure}

Figure~\ref{fig:dutynull} is the theorem measured on paper;
Figure~\ref{fig:ear} was the same theorem measured in an ear, and
Corollary~\ref{cor:hard}(b) is why the two agree: a hard-edged family has
one set of harmonics, and every observer inherits it. A strobe is a third
place to measure it. A camera taking $f$ frames per second is a family
whose picture is a comb, every harmonic equal, and a wheel whose spokes
pass at $r$ per second beats against it: near $r=f$ the wheel is seen
turning at $r-f$, backwards below the frame rate, the wagon-wheel effect
of every western; at $r=f/2$ consecutive frames sit half a spoke apart, a
pooling eye averages them, and what stands still is a wheel with
\emph{twice} the spokes, carried by the spoke profile's second harmonic.
So a wheel whose spokes are half the gap wide cannot show it. Simulated
on sampled frames (script \texttt{wagonwheel.mjs}): the reversal is exact
to $\wagonReversalErr$, the doubled wheel at spoke duty $0.3$ has contrast
$\wagonDoubledAtThirty$ and tracks $|\sin(2\pi d)/2\pi|$ within
$\wagonLawPct\%$, and at duty $\tfrac12$ it is $\wagonDoubledDepth\times$
down; the tripled wheel at $r=f/3$ dies at duties $\tfrac13$ and
$\tfrac23$ ($\wagonTripledDepthLow\times$, $\wagonTripledDepthHigh\times$).
The frames are simulated and the eye is not; the prediction for a person
is that a wheel with fat spokes never shows the doubled still wheel.

Two readings. Forwards, a band as an instrument: find the duty at which a
station dies and you have measured the family's duty to the precision of
the null. Backwards, a warning about what ``the beat'' is: a $16.4{:}8$
pair also carries the recipe $(1,-2)$, one coarse count minus two fine, a
real coefficient with a period of five units and no band anyone can see,
and the first version of this experiment measured a null of \emph{that}
and called it the station. The station is the slow recipe, not the small
one.

\subsection{Which convergents are stations, and which ratios never beat}
\label{sec:stations}

\begin{proposition}[Stations are the convergents with large complete quotients]
\label{prop:stations}
Let a pair have pitch ratio $x=[a_0;a_1,a_2,\dots]$ with convergents
$h_n/k_n$ and complete quotients $x_{n+1}$. The priced merit of the recipe
$(h_n,-k_n)$ is exactly
\[
  \mu_n\;=\;\frac{2\,h_nk_n}{(h_n+k_nx)\,(x_{n+1}k_n+k_{n-1})},
\]
and the convergent is a station, $\mu_n<\tfrac14$, exactly when
$x_{n+1}+k_{n-1}/k_n>8/(1+k_nx/h_n)$, a bound within a few percent of $4$.
A next partial quotient $a_{n+1}\ge5$ guarantees a station and
$a_{n+1}\le2$ forbids one. For the golden ratio $\mu_n\to1/\sqrt5$, for
$\sqrt2$ and the silver ratio it tends to $1/2\sqrt2$: none of those pairs
has a station at any order.
\end{proposition}

\begin{proof}
Put the convergent's error $|h_n-k_nx|=1/(x_{n+1}k_n+k_{n-1})$ from the
lookup into the merit $2hk\,|h-kx|/(h+kx)$; that is the identity. The
threshold rearranges $\mu_n<\tfrac14$. Since $a_{n+1}<x_{n+1}<a_{n+1}+1$
and $0\le k_{n-1}/k_n<1$, a partial quotient of five or more puts the left
side above $5$ and the right side is at most $8$ over $1$-and-a-bit; a
partial quotient of two or less keeps the left side below $4$, and the
small amount by which the right side can fall short of $4$, which is the
convergent's own error, is smaller still for every denominator beyond the
first few, a case check the script does. For the golden ratio every
complete quotient is the golden ratio $\varphi$ itself and
$k_{n-1}/k_n\to1/\varphi$, so $\mu_n\to1/(\varphi+1/\varphi)=1/\sqrt5$; for
$\sqrt2$ the complete quotients are $1+\sqrt2$ and $k_{n-1}/k_n\to\sqrt2-1$.
\end{proof}

The identity is, up to a factor that tends to one, a classical quantity
(Perron's $k_n|k_nx-h_n|$), and $1/\sqrt5$ is Hurwitz's
constant~\cite{Hurwitz1891Annaherung, Khinchin1964Continued}: the theory of
which beats are visible turns out to be the theory of how well a number can
be approximated, priced slightly differently. What is new is the threshold,
which reads a pair's stations off its partial quotients, and the
\emph{deserts}: a ratio whose partial quotients stay small never beats
visibly at any order, however rich the families' harmonics. Measured on
\stationsConvergents\ convergents of \stationsRatios\ random ratios, the
identity agrees with the tool's own scan to $\stationsIdentityErr$, the
limit holds to $\stationsPerronErr$ past the sixth convergent, the
partial-quotient rules hold without exception, the golden ratio's merit
tends to \stationsGoldenMerit\ and the silver ratio's to
\stationsSilverMerit, both above $\tfrac14$ (\texttt{stations.mjs}).
Table~\ref{tab:stations} reads the proposition on five pairs. The desert is
in sound too: two sawtooth tones --- every harmonic present, falling like
$1/n$ --- at $200$ and $410$ hertz, the ratio $2.05$, carry a station line
at ten hertz through a square-law ear that is $\earDesertRatio\times$
stronger than the strongest slow line of two sawtooths at $200$ hertz and
$200$ times the golden ratio (\texttt{ear.mjs}). The prediction for
hearing: golden intervals of harmonic-rich tones sound rough, but they
never \emph{beat}.

\begin{table}[t]
  \centering\small
  % GENERATED by paper/tools/numbers.mjs -- do not edit by hand.
\begin{tabular}{@{}llp{0.52\linewidth}@{}}
\toprule
Pair & Ratio & First convergents; stations in bold \\
\midrule
$16.4{:}8 = 41{:}20$ & 2.0500 & \textbf{2/1}, \textbf{41/20} \\
$15{:}5$, one turned $2^\circ$ & 2.9982 & 2/1, \textbf{3/1}, \textbf{1640/547}, 8203/2736, 26249/8755, \dots \\
$\sqrt2$ & 1.4142 & 1/1, 3/2, 7/5, 17/12, 41/29, \dots \\
$e$ & 2.7183 & 2/1, 3/1, 8/3, 11/4, \textbf{19/7}, \dots \\
$\pi$ & 3.1416 & \textbf{3/1}, \textbf{22/7}, 333/106, \textbf{355/113}, 103993/33102, \dots \\
\bottomrule
\end{tabular}

  \caption{\figlead{Stations, read off the continued fraction.} For each
  pair, the first convergents of its pitch ratio; the ones the threshold
  makes visible stations are in bold. $\pi$ beats at $3$, $22/7$ and
  $355/113$, because its partial quotients $7$, $15$ and $292$ are large;
  $e$ at $19/7$ and later at $193/71$; $\sqrt2$, whose quotients are all
  $2$, at nothing; the $41{:}20$ pair at its octave. Generated from
  \texttt{stations.json}.}
  \label{tab:stations}
\end{table}

\begin{figure}[t]
  \centering
  \includegraphics[width=\linewidth]{selection-stations.png}
  \caption{\figlead{Stations as places.} Two fans of rays from two centres.
  Along the axis between them the local ratio of the two families' rates
  sweeps through the reals, every rational the price admits owns a stretch
  of the axis where its recipe is the slow one, and the marks show where
  each station's bands are: the ladder of convergents laid out as
  geography. A pair whose ratio varies has its stations at places, not at
  settings. Script \texttt{selectionfigs.mjs}.}
  \label{fig:stations}
\end{figure}

\begin{tryit}{a station you can switch off with a marker}
Two rulers, or two drawn combs, at pitches close to $2{:}1$ --- ten and
nineteen to the inch --- show octave bands. Now thicken the coarse comb's
ticks with a marker until each covers half the gap: the bands vanish while
both combs are plainly still there. Thicken them to two thirds and the
bands come back. In \Moire, two layers of lines at pitches $16.4$ and $8$
with the coarse stroke at width $8.2$ do the same thing in the pooled view,
and the \emph{ratio} view shows the station's region going blank.
\end{tryit}
